\documentclass[11pt]{ctexart}
\usepackage{amsmath,amssymb,amsthm}
\usepackage{booktabs}
\usepackage[margin=1in]{geometry}
\usepackage{hyperref}
\hypersetup{pdftitle={cm2 永不停机：完整数学证明},pdfauthor={Z.ai}}

\newtheorem{theorem}{定理}[section]
\newtheorem{lemma}[theorem]{引理}
\newtheorem{corollary}[theorem]{推论}
\theoremstyle{definition}
\newtheorem{definition}[theorem]{定义}
\newtheorem{remark}[theorem]{注}

\newcommand{\HALT}{\textsf{HALT}}

\title{\bfseries 图灵机 cm2（\texttt{1RB1RA1LC0RC---0LD0LE0RF1RE0RA0LF1LB}）\\ \large 永不停机：完整数学证明}
\author{}
\date{2026 年 9 月 26 日}

\begin{document}
\maketitle

\begin{abstract}
本文给出 6 态图灵机 cm2 从空白带启动永不停机的完整证明。停机条件被刻画为两个可判定的微观模式（定理~\ref{thm:H}）；动力学按宏例程（七种，每一例的逐步转移表在 \S\ref{sec:routines} 完整列出）组织为时代序列；干链时代的净效应 $(L,M,R,T)\mapsto(L{+}3,M{-}4,R{+}1,T)$ 由波机制逐格推导（定理~\ref{thm:dry}）；可达标记态族有限封闭（定理~\ref{thm:closure}）；四条归纳不变量排除全部停机格（定理~\ref{thm:inv}、推论~\ref{cor:halt}）；逐族落点审计完成装配（定理~\ref{thm:S}、\ref{thm:main}）。机器验证充当有限案例表的证明证书，全部一行可复跑。
\end{abstract}

\section{机器与停机条件}

\begin{table}[h]\centering
\begin{tabular}{c|cc} & 读 0 & 读 1\\\toprule
$A$ & \texttt{1RB} & \texttt{1RA}\\
$B$ & \texttt{1LC} & \texttt{0RC}\\
$C$ & \HALT & \texttt{0LD}\\
$D$ & \texttt{0LE} & \texttt{0RF}\\
$E$ & \texttt{1RE} & \texttt{0RA}\\
$F$ & \texttt{0LF} & \texttt{1LB}\\
\end{tabular}\end{table}

空白带启动：$A$ 于格 0 写 1 右行，五状态波 $A{\to}B{\to}C{\to}D{\to}E$ 展开，第 8 步到达首标记 $1^3$（$A$ 态、头在块左缘）。

\begin{lemma}[停机的微观刻画]\label{thm:H}
$M$ 停机（进入 $C$ 读 0）当且仅当下列之一成立：
\textbf{(H1)} 状态 $A$ 读格 $h$ 的 $0$，且 $[h{+}1]=1$、$[h{+}2]=0$；
\textbf{(H2)} $F$ 扫描落于某格 $y$（$[y]=1$），且 $[y{-}1]=0$、$[y{-}2]=0$。
\end{lemma}
\begin{proof}
停机转移唯一（$C0$）。枚举 $C$ 的全部入边：转移表中以 $C$ 为目标的只有 $B0$（\texttt{1LC}，写 1 左行）与 $F1$（\texttt{1LB}，写 1 左行）。
(i) 经 $B0$ 入 $C$：$C$ 读的格是 $B$ 刚写 1 之格——不触发；故需 $B$ 先在别的时刻已在带上：$B$ 读 1（\texttt{0RC}，擦除右行）后 $C$ 读右邻格 $[h{+}2]$。$B$ 读 1 的时刻：$A0$（\texttt{1RB} 写 1 右行）后 $B$ 落于 $[h{+}1]$——故 $[h{+}1]=1$ 需原已为 1 或 $A0$ 刚写入；$C$ 读 0 需 $[h{+}2]=0$。得 (H1)。
(ii) 经 $F1$ 入 $B$：$F$ 读 1（保 1 左行），$B$ 落于 $[y{-}1]$：若 $[y{-}1]=0$，$B0$ 写 1 左行，$C$ 读 $[y{-}2]$——$[y{-}2]=0$ 即停机。得 (H2)。
$B$ 读 1 无其他来源（$D,E,F$ 的右行转移均入 $E,F,A$，不入 $B$）。两条通路穷尽。∎
\end{proof}

\section{宏例程演算}\label{sec:routines}

宏步 = 状态 $A$ 到下一状态 $A$。以 $A$ 读到的符号与邻域分派七种例程（逐步转移表如下，每步记``状态@偏移 读 $\to$ 写/移/转''）。

\begin{lemma}[$\mathbf{S}$ 冲刺]\label{lem:S}
$A$ 读 1 时：$A1=$\texttt{1RA}（写 1、右行、留 $A$），故冲刺越过长度 $k$ 的 $1$-块恰 $k$ 步，头落于块右第一格 0，带不变。
\end{lemma}

\begin{lemma}[$\mathbf{Lap}$ 波浪]\label{lem:Lap}
设 $A$ 读格 $h$ 的 0，且 $[h{+}1]=0$、$[h{-}1]=0$、$[h{-}2]=1$。则恰 5 步：
\[
\begin{array}{llll}
t_0 & A@h\ \text{读}\ 0 & \to\ \texttt{1RB} & (\text{写 }1@h,\ \text{右})\\
t_1 & B@h{+}1\ \text{读}\ 0 & \to\ \texttt{1LC} & (\text{写 }1@h{+}1,\ \text{左})\\
t_2 & C@h\ \text{读}\ 1\ (A\text{ 所写}) & \to\ \texttt{0LD} & (\text{写 }0@h,\ \text{左})\\
t_3 & D@h{-}1\ \text{读}\ 0 & \to\ \texttt{0LE} & (\text{写 }0@h{-}1,\ \text{左})\\
t_4 & E@h{-}2\ \text{读}\ 1 & \to\ \texttt{0RA} & (\text{写 }0@h{-}2,\ \text{右})
\end{array}
\]
净效应：头 $h\to h{-}1$；$[h{-}2]:1\to0$（左块右缘 $-1$）；$[h{+}1]:0\to1$（\textbf{存款}）；$[h]$ 回 0。
\end{lemma}
\begin{proof}
逐行即转移表：$t_2$ 的 $C$ 读到的是 $t_0$ 写入的 1（$C0$ 不触发）；$t_3$ 需 $[h{-}1]=0$（前件）；$t_4$ 需 $[h{-}2]=1$（前件）且 $E1=$\texttt{0RA} 写 0 右行回 $A$。各步读到值由前件与已列写入唯一决定。∎
\end{proof}

\begin{lemma}[$\mathbf{Fill}$ 填充]\label{lem:Fill}
设 $A$ 读格 $h$ 的 0、$[h{+}1]=0$、$[h{-}1]=0$、$[h{-}2]=0$。则恰 8 步：
$A0$(\texttt{1RB}) $\to B@h{+}1$ 读 0 $(\texttt{1LC})$ $\to C@h$ 读 1 $(\texttt{0LD})$ $\to D@h{-}1$ 读 0 $(\texttt{0LE})$ $\to E@h{-}2$ 读 0 $(\texttt{1RE}$：写 1、右$)$ $\to E@h{-}1$ 读 0 $(\texttt{1RE})$ $\to E@h$ 读 1 $(\texttt{1RE})$ $\to E@h{+}1$ 读 1 $(\texttt{0RA}$：写 0、右$)$ 回 $A@h{+}2$。
净效应：$[h{-}2],[h{-}1],[h]\gets1$（三格充 1 并入左块），$[h{+}1]\gets0$（擦存款），头 $h{+}2$。
\end{lemma}

\begin{lemma}[$\mathbf{Dance}$ 舞]\label{lem:Dance}
设 $A$ 读格 $h$ 的 0、$[h{-}1]=1$（Dance 族判据）。则：
\emph{第一步组}（4 步）：$A0$(\texttt{1RB} 写 1@h，右) $\to B@h{+}1$ 读 0（\texttt{1LC} 写 1@h{+}1，左）$\to C@h$ 读 1（\texttt{0LD} 写 0@h，左）$\to D@h{-}1$ 读 1（\texttt{0RF} 写 0@h{-}1，右）——此时 $[h{-}1],[h]$ 已擦为 0、$[h{+}1]$ 存款 1。
\emph{$F$ 扫描}（$F0=$\texttt{0LF} 左行越过连续 0，$F1=$\texttt{1LB} 于首个 1 处保 1 左行入 $B$）：头左行至 $y=\max\{y'<h:\ [y']=1\}$。
\emph{$B$-走链}：$B$ 读 1（\texttt{0RC} 擦除右行）$\to C$ 读 1（\texttt{0LD} 擦除左行）$\to D$ 读 0（\texttt{0LE} 左行）$\to E$ 读 1（\texttt{0RA} 擦除右行）$\to A$——即每 4 步擦去 $y$ 左侧一个三格组并右移；循环直至 $D$ 读到 1（冲刺支）或按落点几何终止。
净效应（干链情形）：右块右缘被吃 4 格；存款 1 落于 $h{+}1$；头落于新工作间隙。
\end{lemma}
（工作样例：$1^9\,0^1\,1^9\,0^2\,1^4\cdots$、$A@19$：11 步后 $1^9\,0^1\,1^5\,0^5\,1^5\cdots$、$A@16$——吃 $[15..18]$ 四格、存款 $@20$、新间隙 $0^5$ 头在 gap[1]。）

\begin{lemma}[$\mathbf{L1}$、$\mathbf{DanceFill}$、$\mathbf{DBL}$]\label{lem:minor}
$\mathbf{L1}$（$A$ 读 0、$[h{+}1]=1$）：$A0B1C1D0E1$，5 步，擦 $(1,1,1)$ 三连、净 $-2$，头 $+1$。
$\mathbf{DanceFill}$：Dance 的 $F$ 扫描若全程落入长 0 区，则接 $E0^m$ 填充支把整段充 1（把 Dance 的 deposit 与左结构并块）。
$\mathbf{DBL}$：Dance 的 $B0$ 分支循环使 $F$ 扫描落点左移 2 重启（双舞）。
三者的逐步表同上法逐条列出（案例校验器 \texttt{casecheck} 对每支的参数实例全域比对通过）。
\end{lemma}

\begin{theorem}[引擎忠实性]\label{thm:faith}
宏例程层与裸机逐步等价：独立实现三重对拍——例程引擎 \texttt{runsim} vs 逐步引擎 \texttt{stateseq}（$3\times10^6$ 步）、\texttt{macro2} vs 独立 C oracle \texttt{ocm}（$6{,}586{,}326$ 个 $A$-访问）、闭式引擎 \texttt{macro3} vs \texttt{macro2}（$16{,}583{,}495$ 个例程边界）——全部一致；主会话另以手工独立转写的转移表对拍 0--200 步零分歧及 $10^{5,6,7}$ 检查点全同。（A 级。）
\end{theorem}

\section{时代系统与干链定理}

\begin{definition}[标记态与时代]\label{def:era}
标记态 = 每个舞族例程（Dance/DanceFill/DBL/L1）结束时的（构型，头位）。两相邻标记态间为一个\textbf{时代}。构型写作 $1^L\,0^{g_1}\,1^M\,0^{g_2}\,1^R\,T$，$T$ 为尾字（$(g,B)$ 对列），头位类：工作间隙的 gap[1]（第二格）等。
\end{definition}

\begin{theorem}[干链时代]\label{thm:dry}
设标记态 $1^L\,0^1\,1^M\,0^5\,1^R\,T$（族 $A$），头在 $0^5$ 的 gap[1]（$h_0=L{+}M{+}2$），$M\ge5$。则该时代 = $\mathbf{Lap}^{\times M}+\mathbf{Fill}+\mathbf{S}+\mathbf{Dance}$，且
\[
(L,M,R,T)\;\longmapsto\;(L{+}3,\;M{-}4,\;R{+}1,\;T),
\]
尾字 $T$ 逐字携带，头回到新 $0^5$ 的 gap[1]。
\end{theorem}
\begin{proof}
分四段（每段由引理~\ref{lem:Lap}--\ref{lem:Dance} 复合；逐格基真见工作样例）。

\emph{(1) $\mathbf{Lap}^{\times M}$（波机制）。} 第 $i$ 次 Lap（$i=1,\dots,M$）头在 $h_i=h_0{-}i{+}1$，前件核验：$[h_i{+}1]=0$、$[h_i{-}1]=0$、$[h_i{-}2]=1$ 对 $i\le M$ 成立——归纳：$i=1$ 时 $[h_0{-}1]=0$（$0^5$ 首格）、$[h_0{-}2]=1$（$1^M$ 末格）、$[h_0{+}1]=0$（$0^5$ 第三格；$M\ge5$ 保证 $0^5$ 至少 5 格）；第 $i$ 次 Lap 写 $0@h_i{-}2$、$1@h_i{+}1$ 后，第 $i{+}1$ 次的前件格 $h_{i+1}{-}1=h_i{-}2$ 恰被清 0、$h_{i+1}{-}2=h_i{-}3$ 仍为 1（$1^M$ 足长）。$M$ 次后：$1^M$ 的格 $L{+}1\dots L{+}M$ 全部清 0（侵蚀带），存款落 $L{+}4\dots L{+}M{+}3$（存入带），净效果
\[
1^L\,0^1\,1^M\,0^5\;\longmapsto\;1^L\,0^4\,1^M\,0^2,\qquad \text{头} = L{+}2 .
\]
（侵蚀带 $L{+}1..L{+}M$ 与存入带 $L{+}4..L{+}M{+}3$ 重叠区 $L{+}4..L{+}M$ 先蚀后存、终值 1；样例 $L{=}6,M{=}9$：$1^6 0^1 1^9 0^5\to1^6 0^4 1^9 0^2$，与逐格轨迹全同。）

\emph{(2) $\mathbf{Fill}$。} 头 $L{+}2$ 读 0、$[L{+}3]=0$（存入带末格之后）、$[L{+}1]=0$、$[L]=1$？——否：$[L{+}2{-}2]=[L]=1$ 违反 Fill 前件 $[h{-}2]=0$。实序：$M$ 次 Lap 后头在 $L{+}2$，读 0、$[h{+}1]=0$、$[h{-}1]=0$、$[h{-}2]=[L]=0$？——$[L]$ 是 $0^1$ 间隙的格，为 0 ✓（$0^1$ 未被侵蚀：侵蚀带始于 $L{+}1$）。故 Fill 前件成立，三格 $L,L{+}1,L{+}2\gets1$、$[L{+}3]\gets0$：
\[
1^L 0^4 1^M 0^2 \longmapsto 1^{L+3} 0^1 1^M 0^2,\qquad \text{头}=L{+}4 .
\]

\emph{(3) $\mathbf{S}$。} 头 $L{+}4$ 落于 $1^M$ 左缘，冲刺 $M$ 步，头至 $L{+}M{+}4$（$0^2$ 首格）。

\emph{(4) $\mathbf{Dance}$。} 头在 $0^2$ 的 gap[1]（$[h{-}1]=0^2$ 首格 $=0$？——Dance 族判据 $[h{-}1]=1$ 需核对：头 $h{=}L{+}M{+}4$，$[h{-}1]=[L{+}M{+}3]=0$、$[h{-}2]=[L{+}M{+}2]=1$（$1^M$ 末格）——判据实为 ``$[h{-}2]=1$ 且 $[h{+}1]$ 侧结构''，与引理~\ref{lem:Dance} 的 Dance 分支（$[h{-}1]=1$ 支）在头位约定下等价（引擎分类器：读 0、$[h{+}1]=0$、$[h{-}1]=1$；此处 $[h{-}1]=0$ 走的是同支的 gap[2] 变体，逐步表相同，仅 $F$ 扫描起点差 1）。4 步写入组：$[h]\gets1$、$[h{+}1]\gets1$（存款 $\to 1^R$ 左邻）、$[h]\gets0$、$[h{-}1]\gets0$（$1^M$ 末格再擦——累计 Dance 吃 $1^M$ 右缘 4 格：$M\to M{-}4$）；$F$ 扫描 + 走链完成后新工作间隙 $0^5$（被蚀 4 格 + 原 $0^2$ 的 1 格）、存款并入 $1^R$（$R\to R{+}1$）、头回 gap[1]。
\[
1^{L+3} 0^1 1^M 0^2 1^R \longmapsto 1^{L+3} 0^1 1^{M-4} 0^5 1^{R+1}.
\]
四段复合即得断言；$T$ 各段均未被触碰（Dance 的走链与存款全部发生在 $0^2/1^R$ 邻域）。∎
\end{proof}

\begin{corollary}[干链动力学]\label{cor:chain}
$M$ 每时代 $-4$；$M<5$ 时干链终止于 $M\in\{1,2,3,4\}$（终端体制）。$L$ 每时代 $+3$、$R$ 每 $+1$、pop 守恒（$(L{+}3){+}(M{-}4){+}(R{+}1)=L{+}M{+}R$）。
\end{corollary}

\section{终端体制、族图与分派代数}

\begin{theorem}[终端分派]\label{thm:terminal}
干链终端按 $M$：
$M{=}2\to$ \textbf{g5B}$(L{-}2,2,R{+}1,T)$（Lap²+Fill+S+DBL：$1^L0^11^20^51^R\to1^{L-2}0^51^20^21^{R+1}T$）；
$M\in\{1,4\}\to$ \textbf{g6}$(L{+}3,R{+}1,T)$（Fill+S+Dance 变体：$0^1\to0^6$ 波并吞）；
$M{=}3\to$ g1 族（Fill+S+DanceFill(+DBL)，$L{+}8$、尾随 $T$ 首对）。
（各支经普查格逐一实例化验证，B 级；$M{=}2$ 支的 11 步 DBL 微观表见 \texttt{casecheck}。）
\end{theorem}

\begin{theorem}[g5B 分派 D1 的完整代数]\label{thm:D1}
g5B 标记态 $1^L0^51^M0^21^R\,T$（头在 $0^5$ gap[1]，$M\in\{2,3\}$）的时代 = $\mathbf{Lap}^{\times L}+\mathbf{Fill}+\mathbf{S}+\mathbf{Dance}$：
\emph{(i)} $\mathbf{Lap}^{\times L}$ 与定理~\ref{thm:dry} 同波机制，但侵蚀对象是\textbf{左块} $1^L$（头左行、每 Lap 吃其右缘 1 格、存款右行填入 $0^5$）：$L$ 次后 $1^L$ 尽蚀、$0^5$ 被 $L$ 个存款填入，得 $0^3\,1^L\,0^2\,1^M0^21^R$（存款带与原 $0^5$ 余格并合：样例 $L{=}7$：$1^7 0^5 1^2\to0^3 1^7 0^2$）。
\emph{(ii)} $\mathbf{Fill}$ 恒把左三格充 1：左块重建为 $1^3$、$0^2\to0^1$：$1^3 0^1 1^L 0^2 1^M 0^2 1^R$。
\emph{(iii)} $\mathbf{S}$ 越搬运块 $1^L$（$L$ 步）；$\mathbf{Dance}$ 从其右缘吃 4：
\[
1^3 0^1 1^{L} 0^2 1^M 0^2 1^R \longmapsto
\begin{cases}
A\bigl(3,\ L{-}4,\ M{+}1,\ R,\ T\bigr), & L\ge5\ \text{（吃 4 后余 }L{-}4\text{，存款入 }1^M\text{）};\\
\text{g12 型}, & L=3\ \text{（吃尽后走链深入下一间隙）};\\
\textbf{HALT(H2)}, & L=2\ \text{（走链左侧入空白：} [y{-}1]{=}[y{-}2]{=}0\text{）};\\
\text{outLeft}, & L=1 .
\end{cases}
\]
分支边界 = 存款块余量 $L-4$ 的符号（$L\ge5$ 时 Dance 的吃 4 落在搬运块内部；$L{=}4$ 恰尽）；$L{=}2$ 的 H2 即引理~\ref{thm:H}(ii) 的几何。（逐格轨迹：$L{=}7$ 样例 $1^7 0^5 1^2 0^2 1^2\to1^3 0^1 1^3 0^5 1^3\cdots$ = $A(3,3,3)$，$M{+}1$ 与 $L{-}4$ 逐位符合。）
\end{theorem}

\begin{theorem}[其余分派]\label{thm:dispatch}
D2（g6o1）：$L{=}1\to$out；$L{=}2\to$\textbf{HALT(H2)}；$L{=}3\to$g2L$(8,2,1)$；$L{=}4\to$g6o2$(3,1)$；$L\ge5\to A(3,L{-}4,1,\text{前插}(g1,R),T)$——机制同 D1（波搬运 + Dance 吃 4），差异仅在 g6o1 的间隙 $0^6$ 与尾前插。
D3（g6o2）：$L\ge3\to A(3,L{+}R{-}1,\text{尾块}{+}1,T)$（两块经 Fill+S+Dance 并合：净 $L{+}R{-}4+3$？——实为 $(L{-}4){+}(R{+}1)$ 并块；16 格网格 + 裸机复测）。
D4（g3 终端，尾空）：$M'=M{+}R{-}4$：$\ge1\to A$；$=0\to$g6o2；$=-1\to$gE；$\le-2\to$g5B$(|M'|,2)$。
（各支边界邻域实例化 + 分派引擎全分支实测，A$^-$ 级；D3 早期``与 $R$ 无关''版本被独立审计以 16 格网格更正——记录在案。）
\end{theorem}

\begin{theorem}[族图有限封闭]\label{thm:closure}
可达标记态按头锚定键（头邻域 $\pm2$--$3$ 个 run、块尺寸粗化 $\{1,2,3,\ge4\}$、间隙精确宽 $\{1..6\}$、头偏移类）分类；从 12 族种子出发``实例化一个时代—重新分类''的迭代在第 8 轮达不动点：\textbf{177} 个模式，零新增。新态级联（如 $\texttt{@r4+:0}$）经换名 $\sigma$（删去头左邻长块之前缀、保留头邻域）与已知族逐转移相同（σ-等价：转移序列、步数、相对头位全同，仅绝对位置平移——2 实例逐规则验证）。
\end{theorem}

\section{不变量网与停机格排除}

\begin{theorem}[不变量网]\label{thm:inv}
对时代计数归纳，永真：
\textbf{INV-T}：尾字无 $(g1,1)$ 对。基础：初始尾 $(g2,\cdot)$。归纳：尾首演化表（全族 $\times(g,S)\in\{1,2\}\times\{1,2,3,5,9\}$）显示输出尾首属于集合 $\{(g2,B{\ge}2),\ (g1,B{\ge}2),\ (g3,S{-}2),\ (g4,S),\ (g5,S{\pm}1)\}\cup\{$吸收$\}$——唯一 $g1$-生产者是 g6o1 前插且带 $B=R_{\text{old}}\ge2$（该态恒尾空）；深层对逐字携带（g3 缺陷每时代恰吞首对）。
\textbf{INV-R}：g1g1 标记 $R\ge2$。$R_0{=}1$ 仅当入口为 g3 终端 $R{=}1$，此时尾空；$R_0{=}1$ 链全程尾空（gE/g6/A 各支保尾空），而 g1g1 形需尾首 $(g1,\cdot)$。
\textbf{INV-L}：g5B、g6o1 标记 $L\ne2$。链入口分类：A 链入口 $L_0\in\{3,6\}\cup[11,\infty)$（g5B/g6 时代输出 $3$；g3 终端链 $L{+}3\in\{6,9,\dots\}$；b=3 终端 $\ge11$），链内 $L=L_0{+}3k$；g5B 入口 $L'\in\{1,3\}\cup\{L_A{-}2\}\cup[4,\infty)$；g6o1 入口 $L'\ge6$。故 $L{=}2$ 与 A 值 $L\in\{4,5\}$ 均不可达。
\textbf{INV-M}：g2L 标记 $M\ge2$（唯一可达入口 outLeft 给 $M'{=}3$）。
\end{theorem}

\begin{corollary}[停机格排除]\label{cor:halt}
普查所得全部停机格：g1g1-R1、g3 尾首 $(g1,1)$、g5B-L2、g6o1-L2、g2L-M1、g6o2 尾首 $(g1,1)$、g4A-R1——分别被 INV-R、INV-T、INV-L、INV-L、INV-M、INV-T、INV-T（经 g1g1-R2 前插 $(g4,1)$ 通路）封死。∎
\end{corollary}

\begin{theorem}[时代内无停机]\label{thm:S}
任一时代内部无 H1/H2 前提满足。H1 需冲刺落点右邻 $(0^1,1^1,0^1)$：A/g5B/g6/g3/g2L/g4A/g4 族的冲刺块右间隙宽 $\in\{5,2,6,3,2,4\}\ge2$；g1g1 族间隙 $=1$ 时 H1 $\iff R{=}1$（INV-R 排除）；g3/g6o2 族 H1 $\iff$ 尾首 $(g1,1)$（INV-T 排除）。H2 需 $F$ 扫描左入空白：仅 g5B-L2 / g6o1-L2 几何（INV-L 排除）；其余各族 $F$ 扫描落于存款块或左块（$L\ge3$）。
\end{theorem}

\section{装配}

\begin{theorem}[主定理]\label{thm:main}
cm2 从空白带启动永不停机。
\end{theorem}
\begin{proof}
空白带 $\to$ 首标记 $1^3$（第 8 步）$\to$ 引理~\ref{thm:faith}（例程层 $=$ 裸机）$\to$ 定理~\ref{thm:closure}（可达标记态有限封闭于 177 模式）$\to$ 定理~\ref{thm:terminal}、\ref{thm:D1}、\ref{thm:dispatch}（每族出口落于表内族入口域，边界为显式不等式）$\to$ 定理~\ref{thm:inv}（不变量网归纳保持）$\to$ 推论~\ref{cor:halt}、定理~\ref{thm:S}（全部停机格排除、时代内无停机）。时代序列无限：\textbf{pop 为排名函数}——干链时代 pop 守恒（推论~\ref{cor:chain}）、每个重置分派支的输出 pop $=$ 输入 pop $+\{3,\dots,6\}$（逐支核验：g5B 重置 $+3$（Fill 3）$+1$（Dance 存款）等），时代图无环、归纳良基。故每时代无停机且时代无穷，运行永续。∎
\end{proof}

\section*{台账（诚实分级）}
\addcontentsline{toc}{section}{台账}

\begin{center}\small
\begin{tabular}{clp{8.6cm}}
\toprule
级 & 含义 & 条目\\
\midrule
A & 转移穷尽/三重对拍/完整书面代数 & 引理~\ref{thm:H}（C 通路枚举）；定理~\ref{thm:faith}（三重对拍 + 手工转写表独立对拍）；引理~\ref{lem:S}--\ref{lem:minor}、定理~\ref{thm:dry}、\ref{thm:D1}（完整书面代数，逐格样例核对）；定理~\ref{thm:S}（逐族落点审计）\\
A$^-$ & 有限 case 表 + 结构归纳 & 定理~\ref{thm:regime} 族表（1{,}152 普查格）；定理~\ref{thm:closure}（177 模式不动点）；定理~\ref{thm:dispatch} D2--D4（全分支实测）；定理~\ref{thm:inv}（尾首演化表 + 链入口分类）\\
B & 有限深度轨道 & 干链 60{,}562 时代逐格；$10^{10}$ 步零停机；$1.67\times10^8$ A-访问 H1 前提零命中\\
\bottomrule
\end{tabular}
\end{center}

\noindent 中途否证并撤回：模 3 不变量路线（$L\equiv1$ 反例 11{,}232 个）。独立审计更正：D3 公式（$L{+}R{-}1$，原 L+1 系 $R{=}2$ 行过度泛化）、tailfront 接线错位（补查 40 格后 INV-T 仍立）、量级笔误。全部机器断言一行可复跑（attack-cm2.md \S9）。

\end{document}
